\begin{theorem}
    Every prime element is irreducible.
\end{theorem}
\begin{proof}
    Let \( R \) be an integral domain, and let \( p \) be a prime element in \( R \). By definition, \( p \) is prime if whenever \( p \) divides a product \( ab \), it must divide at least one of \( a \) or \( b \), i.e.,  
\[
p \mid ab \Rightarrow p \mid a \text{ or } p \mid b.
\]  
We want to show that \( p \) is irreducible, meaning that if \( p = ab \) for some \( a, b \in R \), then either \( a \) or \( b \) is a unit in \( R \).  

Since \( p = ab \), we apply the definition of primeness to conclude that \( p \mid a \) or \( p \mid b \). Without loss of generality, suppose \( p \mid a \). Then, there exists some \( c \in R \) such that \( a = pc \). Substituting this into \( p = ab \), we get:  
\[
p = (pc)b = p(cb).
\]  
Since \( R \) is an integral domain, we can cancel \( p \), yielding \( 1 = cb \), which shows that \( b \) is a unit. Thus, \( p \) is irreducible.
\end{proof}
But every irreducible element is not prime.
\section{Unique Factorization Domain}
\begin{definition}
    An integral domain \( U \) is a \textit{unique factorization domain (UFD)} if every non-zero, non-unit element of \( U \) can be written as a product of irreducible elements, and this factorization is unique up to order and multiplication by units. 
    Moreover, in a UFD, every irreducible element is also a prime element.
    \end{definition}
    Let \( p \) be an element of a UFD \( U \). Suppose \( p \) has two factorizations into irreducibles:
    \[
    p = a_1 \cdots a_n = q_1 \cdots q_m.
    \]
    Since \( U \) is a UFD, uniqueness of factorization implies that each \( a_i \) is a unit multiple of some \( q_j \), meaning there exist units \( u_i \) such that
    \[
    a_1 = u_1 q_1, \quad a_2 = u_2 q_2, \quad \dots, \quad a_n = u_n q_n.
    \]
    From the standard results:
    \[
    \text{ED} \Rightarrow \text{PID} \Rightarrow \text{UFD}.
    \]
    For example:
    \begin{itemize}
        \item \( K[x,y] \) (a polynomial ring over a field \( K \)) is a UFD but not a PID.
        \item \( \mathbb{Z} \) is an ED, and hence also a PID and a UFD.
    \end{itemize}
    To see that \( (2,3) = \mathbb{Z} \), observe that:
    \[
    (2,3) = \{ 2k + 3l \mid k, l \in \mathbb{Z} \}.
    \]
    Since \( 1 \) can be written as \( 1 = 2(2) + 3(-1) \), it follows that \( (2,3) = (1) = \mathbb{Z} \).